\originalchapter{匹配与覆盖}\label{ch:matching}
\sourceof{18.315 L7、L24; 18.433 L1--3; Tutte 极大超图证明为编者另证}
\originalsection{增广路}
\begin{definition}
匹配是两两不共端点的边集 $M$. 顶点若关联某条匹配边就称为已匹配, 否则为自由顶点. 匹配覆盖全部顶点时为完美匹配. 匹配数 $\nu(G)$ 是匹配的最大边数. 顶点覆盖是与每条边都至少有一端相交的顶点集, 最小大小记 $\tau(G)$.
\end{definition}
极大匹配不能再加入一条边, 但可能并非最大匹配. 在四顶点路 $1,2,3,4$ 上, 只有中间边的匹配是极大匹配, 而两端边组成大小2的最大匹配. 要改进前者, 必须先删除一条匹配边, 再加入两条边.
\begin{definition}
相对于 $M$, 交替路的边交替属于和不属于 $M$. 增广路是两个端点都是自由顶点、首尾边都不在 $M$ 的交替路. 对增广路翻转边的归属, 得到边数增加1的新匹配.
\end{definition}
\begin{theorem}[Berge 判据]\label{thm:berge}
匹配 $M$ 最大当且仅当不存在相对于它的增广路.
\end{theorem}
\begin{proof}
有增广路就能翻转而增加匹配, 所以非最大. 反向若 $M'$ 更大, 对称差 $M\mathbin\triangle M'$ 的每个顶点度至多2, 因而分支为交替圈、交替路或孤立点. 圈中两种边同样多, 偶长路也同样多. 因总计 $|M'|>|M|$, 至少一路中 $M'$ 边比 $M$ 边多1. 首尾为 $M'$ 边, 端点在 $M$ 中自由, 这就是增广路.
\end{proof}
\begin{center}
\begin{tikzpicture}[x=1.3cm]
\foreach \i/\lab in {0/a,1/u,2/b,3/v,4/c,5/w}\node[vertex](\lab)at(\i,0){$\lab$};
\foreach \u/\v in {a/u,b/v,c/w}\draw[edge,dashed](\u)--(\v);
\foreach \u/\v in {u/b,v/c}\draw[edge,double distance=1.5pt](\u)--(\v);
\end{tikzpicture}
\end{center}
双线为当前匹配, 虚线为非匹配边. 端点 $a,w$ 自由; 翻转整条路后匹配从2条变为3条. 中间顶点仍恰被一条匹配边覆盖.
\originalsection{Hall 定理}
\begin{theorem}[Hall]\label{thm:hall}
二部图 $G=(A\sqcup B,E)$ 存在覆盖全部 $A$ 的匹配, 当且仅当每个 $S\subseteq A$ 都满足 $|N(S)|\ge|S|$.
\end{theorem}
\begin{proof}
必要性: $S$ 中各顶点的匹配伙伴不同且都在 $N(S)$ 中. 充分性对 $|A|$ 归纳. 空侧无须匹配, 单顶点由条件有邻居. 若每个非空真子集 $S\subsetneq A$ 都满足严格不等式, 任取边 $ab$ 并删除其两端; 对剩余侧的非空 $S$, 原邻域至少 $|S|+1$, 删除 $b$ 后仍不少于 $|S|$. 归纳得到剩余匹配, 加回 $ab$ 即可.

若有非空真子集 $S$ 满足 $|N(S)|=|S|$, 在 $S\sqcup N(S)$ 上条件仍成立, 先归纳匹配这部分. 对余图, 任取 $T\subseteq A\setminus S$, 原条件用于 $S\cup T$ 得
\[
|N(S)\cup N(T)|\ge |S|+|T|,
\]
减去 $|N(S)|=|S|$ 得 $|N(T)\setminus N(S)|\ge|T|$. 余图也满足条件, 再匹配后合并两个匹配.
\end{proof}
Hall 条件一次检查所有子集, 并不是说逐个顶点有邻居就够. 三个工人可选的工作分别为 $\{1,2\}$、$\{1,2\}$、$\{1,2\}$, 每人有两选择, 但三人共同邻域只有两项, 不可能全部安排.
\begin{corollary}\label{cor:regularmatching}
有限 $k$ 正则二部多重图在 $k\ge1$ 时有完美匹配.
\end{corollary}
\begin{proof}
两侧度数和相等, 故 $k|A|=k|B|$, 两侧等大. 对 $S\subseteq A$, 从 $S$ 发出的 $k|S|$ 条边都落入 $N(S)$, 而那里至多承接 $k|N(S)|$ 条边. 所以满足 Hall 条件. 定理证明同样适用于多重边, 因为邻域与顶点匹配问题不因重复边改变.
\end{proof}
\originalsection{Kőnig 定理与算法证书}
\begin{theorem}[Kőnig]\label{thm:konig}
有限二部图满足 $\nu(G)=\tau(G)$.
\end{theorem}
\begin{proof}
任意覆盖至少要从每条匹配边选一个端点, 所以 $\tau\ge\nu$. 取最大匹配 $M$, 从 $A$ 侧的所有自由顶点出发, 沿“非匹配边由 $A$ 到 $B$, 匹配边由 $B$ 到 $A$”访问可达顶点. 记可达部分为 $Z_A,Z_B$. 由 Berge, $Z_B$ 没有自由顶点, 否则找到增广路. 每个 $Z_B$ 顶点的匹配伙伴在 $Z_A$ 中; 每个已匹配的 $Z_A$ 顶点的伙伴也在 $Z_B$ 中, 因为到达它的最后一步必须沿该匹配边.

令 $C=(A\setminus Z_A)\cup Z_B$. 若边 $ab$ 没被覆盖, 就有 $a\in Z_A,b\notin Z_B$. 非匹配边会使 $b$ 可达; 匹配边则由伙伴性质使 $b$ 可达, 都矛盾. 因而 $C$ 是覆盖. 每条匹配边的两端同在可达部分或同在不可达部分, 在 $C$ 中恰选一端; 自由顶点均不在 $C$ 中. 所以 $|C|=|M|$.
\end{proof}
这同时给出最大匹配的可检查证书: 提交一个匹配和同样大小的覆盖, 两者分别给出下界与上界. 二部匹配算法从空匹配开始反复搜索增广路, 至多增广 $\min(|A|,|B|)$ 次. 一次搜索扫描顶点与边, 时间 $O(n+m)$, 因而朴素总时间 $O(n(n+m))$. 非二部图中搜索树可能遇到奇圈, 需用 Edmonds 花收缩算法; 本册不把二部搜索直接套到一般图.
\originalsection{一般图的完美匹配}
\begin{theorem}[Tutte]\label{thm:tutte-matching}
有限简单图 $G$ 有完美匹配当且仅当每个 $S\subseteq V$ 都满足 $o(G-S)\le|S|$, 其中 $o(F)$ 为 $F$ 的奇阶分支数.
\end{theorem}
\begin{proof}
若有完美匹配, $G-S$ 的每个奇阶分支至少有一个顶点要与 $S$ 中顶点匹配; 各分支占用不同顶点, 得必要性. 空图也符合定理.

证明充分性. 条件在 $S=\varnothing$ 时保证各分支阶数为偶数, 故总阶数为偶数. 假设没有完美匹配, 在同一顶点集上加边, 取边数极大的无完美匹配超图 $H$. 对每条缺失边 $e$, 图 $H+e$ 都有使用 $e$ 的完美匹配. 令 $U$ 为 $H$ 中与所有其他点相邻的顶点集.

先证 $H-U$ 的每个分支为团. 否则非完全连通分支含诱导路 $x,y,z$, $xz$ 不存在. 因 $y\notin U$, 有顶点 $w$ 不邻接 $y$; 它不同于 $x,y,z$. 取 $H+xz$ 的完美匹配 $M_1$ 和 $H+yw$ 的完美匹配 $M_2$. 两匹配对称差分解成交替偶圈. 如果新增边位于不同圈, 在 $M_1$ 中翻转包含 $xz$ 的圈, 就得到 $H$ 内的完美匹配, 矛盾.

所以两条新增边同在一交替圈 $C$. 删去它们后, 得两条偶长交替路, 每条连接 $\{x,z\}$ 中一点与 $\{y,w\}$ 中一点. 在每条路上取 $M_1$ 边, 只留下来自 $\{x,z\}$ 的端点; 取 $M_2$ 边, 则只留下来自 $\{y,w\}$ 的端点. 若两路分别连接 $x$ 到 $y$、$z$ 到 $w$, 第一条取 $M_2$, 第二条取 $M_1$, 再加原边 $yz$. 若分别连接 $x$ 到 $w$、$z$ 到 $y$, 第一条取 $M_1$, 第二条取 $M_2$, 再加原边 $xy$. 都将圈上各点用 $H$ 的边匹配, 圈外沿用 $M_1$, 矛盾. 结构结论成立.

设 $H-U$ 的奇阶团数为 $q$. 若 $q\le|U|$, 每个奇团取一点与不同的 $U$ 顶点匹配, 各团剩余偶数点在内部配对. $U$ 剩余点数 $|U|-q$ 为偶数, 因总阶数的奇偶性等于 $|U|+q$; $U$ 自身也是团, 所以能内部配对. 又得到完美匹配, 矛盾. 故 $o(H-U)>|U|$. 由于 $G\subseteq H$, 每个 $H-U$ 的奇分支在 $G-U$ 中细分后至少贡献一个奇分支, 所以 $o(G-U)\ge o(H-U)>|U|$, 与原条件矛盾.
\end{proof}
Tutte 条件揭示了真正障碍是“太多奇分支争夺太少外部匹配伙伴”. 例如六顶点星 $K_{1,5}$, 删中心后有5个奇分支, 大于1, 所以偶数阶并不足以保证完美匹配.
\originalsection{习题与解答}
\begin{exercise}\label{ex:match1}
用 Hall 定理证明: 任意 $n\times n$ 非负矩阵若每行、每列和均为1, 则能选出一个置换 $\pi$ 使每个 $a_{i,\pi(i)}>0$.
\end{exercise}
\begin{exercise}\label{ex:match2}
证明任意极大匹配的大小至少是最大匹配的一半.
\end{exercise}
习题\ref{ex:match1}的解答.
\begin{solution}
把行、列作为两侧顶点, 正元素位置连边. 对行集 $S$, 它们总和为 $|S|$, 所有贡献位于 $N(S)$ 中的列, 每列总和为1, 因而 $|S|\le|N(S)|$. 完美匹配对应所求置换.
\end{solution}
习题\ref{ex:match2}的解答.
\begin{solution}
极大匹配 $M$ 的所有端点构成覆盖: 若存在两端都自由的边就还能加入. 因而每条最大匹配边至少占用其中一个端点, 且不同最大匹配边占用不同端点. 得 $\nu\le2|M|$. 四顶点路上的中间边例子达到等号.
\end{solution}
